% IEEE-style reference entries keyed for \cite; build_report.py orders them by first citation.
\bibitem{idf2022} H. Sun \emph{et al.}, ``IDF Diabetes Atlas: Global, regional and country-level diabetes prevalence estimates for 2021 and projections for 2045,'' \emph{Diabetes Res. Clin. Pract.}, vol. 183, Art. no. 109119, 2022, doi: 10.1016/j.diabres.2021.109119.
\bibitem{cdc2015data} Centers for Disease Control and Prevention, ``2015 BRFSS survey data and documentation,'' Behavioral Risk Factor Surveillance System, Aug. 2016. [Online]. Available: https://www.cdc.gov/brfss/annual\_data/annual\_2015.html
\bibitem{cdc2015codebook} Centers for Disease Control and Prevention, ``Behavioral Risk Factor Surveillance System 2015 codebook report: Land-line and cell-phone data,'' Aug. 23, 2016. [Online]. Available: https://www.cdc.gov/brfss/annual\_data/2015/pdf/codebook15\_llcp.pdf
\bibitem{idfafrica2021} International Diabetes Federation Africa Region, ``24m adults are living with diabetes in Africa,'' IDF, 2021. [Online]. Available: https://idf.org/?p=948
\bibitem{mohamed2018} S. F. Mohamed \emph{et al.}, ``Prevalence and factors associated with pre-diabetes and diabetes mellitus in Kenya: Results from a national survey,'' \emph{BMC Public Health}, vol. 18, suppl. 3, Art. no. 1215, 2018, doi: 10.1186/s12889-018-6053-x.
\bibitem{pierannunzi2013} C. Pierannunzi, S. S. Hu, and L. Balluz, ``A systematic review of publications assessing reliability and validity of the Behavioral Risk Factor Surveillance System (BRFSS), 2004--2011,'' \emph{BMC Med. Res. Methodol.}, vol. 13, Art. no. 49, 2013, doi: 10.1186/1471-2288-13-49.
\bibitem{xie2019} Z. Xie, O. Nikolayeva, J. Luo, and D. Li, ``Building risk prediction models for type 2 diabetes using machine learning techniques,'' \emph{Prev. Chronic Dis.}, vol. 16, Art. no. 190109, 2019, doi: 10.5888/pcd16.190109.
\bibitem{majcherek2025} D. Majcherek, A. Ciesielski, and P. Sobczak, ``AI-driven analysis of diabetes risk determinants in U.S. adults: Exploring disease prevalence and health factors,'' \emph{PLOS ONE}, vol. 20, no. 9, Art. no. e0328655, 2025, doi: 10.1371/journal.pone.0328655.
\bibitem{zafar2025} M. M. Zafar, Z. A. Khan, N. Javaid, M. Aslam, and N. Alrajeh, ``From data to diagnosis: A novel deep learning model for early and accurate diabetes prediction,'' \emph{Healthcare}, vol. 13, no. 17, Art. no. 2138, 2025, doi: 10.3390/healthcare13172138.
\bibitem{chowdhury2026} S. S. Chowdhury, N. Jahan, S. Mukhopadhyay, S. Fang, and V. R. Ramakrishnan, ``The ceiling is in the channel: Auditing learner gaps and measurement frontiers in clinical prediction,'' arXiv:2609.01909, Sep. 2026 (preprint).
\bibitem{christodoulou2019} E. Christodoulou, J. Ma, G. S. Collins, E. W. Steyerberg, J. Y. Verbakel, and B. Van Calster, ``A systematic review shows no performance benefit of machine learning over logistic regression for clinical prediction models,'' \emph{J. Clin. Epidemiol.}, vol. 110, pp. 12--22, 2019, doi: 10.1016/j.jclinepi.2019.02.004.
\bibitem{kopitar2020} L. Kopitar, P. Kocbek, L. Cilar, A. Sheikh, and G. \v{S}tiglic, ``Early detection of type 2 diabetes mellitus using machine learning-based prediction models,'' \emph{Sci. Rep.}, vol. 10, Art. no. 11981, 2020, doi: 10.1038/s41598-020-68771-z.
\bibitem{grinsztajn2022} L. Grinsztajn, E. Oyallon, and G. Varoquaux, ``Why do tree-based models still outperform deep learning on typical tabular data?'' in \emph{Adv. Neural Inf. Process. Syst. (NeurIPS), Datasets and Benchmarks Track}, 2022. [Online]. Available: https://arxiv.org/abs/2207.08815
\bibitem{shwartzziv2022} R. Shwartz-Ziv and A. Armon, ``Tabular data: Deep learning is not all you need,'' \emph{Inf. Fusion}, vol. 81, pp. 84--90, 2022, doi: 10.1016/j.inffus.2021.11.011.
\bibitem{friedman2001} J. H. Friedman, ``Greedy function approximation: A gradient boosting machine,'' \emph{Ann. Statist.}, vol. 29, no. 5, pp. 1189--1232, 2001, doi: 10.1214/aos/1013203451.
\bibitem{breiman2001} L. Breiman, ``Random forests,'' \emph{Mach. Learn.}, vol. 45, no. 1, pp. 5--32, 2001, doi: 10.1023/A:1010933404324.
\bibitem{cheng2016} H.-T. Cheng \emph{et al.}, ``Wide \& deep learning for recommender systems,'' in \emph{Proc. 1st Workshop Deep Learn. Recommender Syst. (DLRS)}, 2016, pp. 7--10, doi: 10.1145/2988450.2988454.
\bibitem{saito2015} T. Saito and M. Rehmsmeier, ``The precision-recall plot is more informative than the ROC plot when evaluating binary classifiers on imbalanced datasets,'' \emph{PLOS ONE}, vol. 10, no. 3, Art. no. e0118432, 2015, doi: 10.1371/journal.pone.0118432.
\bibitem{collins2024} G. S. Collins \emph{et al.}, ``TRIPOD+AI statement: Updated guidance for reporting clinical prediction models that use regression or machine learning methods,'' \emph{BMJ}, vol. 385, Art. no. e078378, 2024, doi: 10.1136/bmj-2023-078378.
\bibitem{lundberg2017} S. M. Lundberg and S.-I. Lee, ``A unified approach to interpreting model predictions,'' in \emph{Adv. Neural Inf. Process. Syst. 30 (NIPS)}, 2017. [Online]. Available: https://papers.nips.cc/paper/7062-a-unified-approach-to-interpreting-model-predictions
\bibitem{strobl2007} C. Strobl, A.-L. Boulesteix, A. Zeileis, and T. Hothorn, ``Bias in random forest variable importance measures: Illustrations, sources and a solution,'' \emph{BMC Bioinformatics}, vol. 8, Art. no. 25, 2007, doi: 10.1186/1471-2105-8-25.
\bibitem{obermeyer2019} Z. Obermeyer, B. Powers, C. Vogeli, and S. Mullainathan, ``Dissecting racial bias in an algorithm used to manage the health of populations,'' \emph{Science}, vol. 366, no. 6464, pp. 447--453, 2019, doi: 10.1126/science.aax2342.
\bibitem{pedregosa2011} F. Pedregosa \emph{et al.}, ``Scikit-learn: Machine learning in Python,'' \emph{J. Mach. Learn. Res.}, vol. 12, pp. 2825--2830, 2011.
\bibitem{abadi2016} M. Abadi \emph{et al.}, ``TensorFlow: A system for large-scale machine learning,'' in \emph{Proc. 12th USENIX Symp. Operating Syst. Design Implementation (OSDI)}, 2016, pp. 265--283.
